\section{Erdős–Rényi Graphs}
\label{sec:random-graphs}

We study the \MinCutPath{} problem in Erdős–Rényi random graphs $G(n, p)$, in which each of the possible edges on $n$ vertices is present independently with probability $p$.
Throughout this section, $\log$ denotes the natural logarithm.
An event holds \emph{with high probability} if its probability tends to one as $n \to \infty$.

Dense random graphs, i.e., those with a constant edge probability $p$, have diameter two with high probability~\cite{Bollobas2001}.

\begin{theorem}[\cite{Bollobas2001}]
\label{thm:random-diameter-two}
Let $G(n, p)$ be an Erdős–Rényi random graph with $n$ vertices and edge probability $p = 1/\alpha$, where $\alpha > 1$ is a constant.
Then the probability that $G(n, p)$ has diameter two converges to 1 as $n \to \infty$:
\[
\lim_{n \to \infty} \mathbb{P}\bigl(\diam(G(n, 1/\alpha)) = 2\bigr) = 1.
\]
\end{theorem}

By Theorem~\ref{thm:diameter-two} and Remark~\ref{rem:algorithm}, a minimum cut-path in a graph of diameter two has $c(u, v) + d(u, v) - 1$ edges and can be computed in polynomial time.
Consequently, for dense random graphs, the \MinCutPath{} problem can be solved exactly in polynomial time on almost all inputs.

Consider now sparse Erdős–Rényi random graphs $G(n, p)$ with $p = \alpha \log n / n$, where $\alpha > 1$ is a constant.
For such graphs, we use two known results, on the diameter and on the edge connectivity~\cite{Bollobas2001,ChungLu2001,Frieze2016}.
We also bound the degrees (Lemma~\ref{lem:degree-bounds}) and the values $c(u,v)$ (Lemma~\ref{lem:connectivity-bounds}).

\begin{theorem}[Diameter~\cite{Bollobas2001,ChungLu2001}]
\label{thm:random-diameter}
Let $G(n, p)$ be an Erdős–Rényi random graph with $p = \alpha \log n / n$, where $\alpha > 1$.
Then
\[
\lim_{n\to\infty} \mathbb{P}\left( \diam(G) \leq \frac{\log n}{\log \log n} \right) = 1.
\]
\end{theorem}

\begin{theorem}[Connectivity~\cite{Frieze2016}]
\label{thm:random-connectivity}
Let $G(n, p)$ be an Erdős–Rényi random graph with $p = \alpha \log n / n$, where $\alpha > 1$.
Then
\[
\lim_{n\to\infty} \mathbb{P}\bigl( G \text{ is } k\text{-edge-connected} \bigr) = 1,
\]
where $k = \delta(G)$ denotes the minimum degree of $G$.
\end{theorem}

\begin{lemma}[Degree Bounds]
\label{lem:degree-bounds}
Let $G(n, p)$ be an Erdős–Rényi random graph with $p = \alpha \log n / n$, where $\alpha > 1$ is a fixed constant.
Then there exist constants $0 < \beta_1 < \beta_2$, depending only on $\alpha$, such that
\[
\lim_{n\to\infty} \mathbb{P}\Bigl( \forall v \in V(G): \ \beta_1 \log n \leq \deg(v) \leq \beta_2 \log n \Bigr) = 1.
\]
\end{lemma}

\begin{proof}
The degree of a fixed vertex $v$ has the binomial distribution with mean $\mu = (n-1)p = (1-o(1))\,\alpha \log n$.
For $a > 0$, let $h(a) = a \log a - a + 1$.
By the Chernoff bounds (see, e.g.,~\cite{Frieze2016}),
\[
\mathbb{P}\bigl(\deg(v) \leq a\mu\bigr) \leq \exp(-\mu h(a)) \ \text{ for } 0 < a < 1
\qquad\text{and}\qquad
\mathbb{P}\bigl(\deg(v) \geq a\mu\bigr) \leq \exp(-\mu h(a)) \ \text{ for } a > 1.
\]
Since $h(a) \to 1$ as $a \to 0$, $h(a) \to \infty$ as $a \to \infty$, and $\alpha > 1$, we can fix constants $0 < a_1 < 1 < a_2$ such that $\alpha h(a_1) > 1$ and $\alpha h(a_2) > 1$.
For these constants, both probabilities above are $o(1/n)$, because $\exp(-\mu h(a_i)) = n^{-(1-o(1))\,\alpha h(a_i)}$ and $\alpha h(a_i) > 1$ for $i = 1, 2$.
The union bound over all $n$ vertices shows that, with high probability, every degree lies between $a_1 \mu$ and $a_2 \mu$.
Since $\mu = (1 - 1/n)\,\alpha \log n$, we have $\alpha \log n / 2 \leq \mu \leq \alpha \log n$ for $n \geq 2$.
The claim follows, for instance, with $\beta_1 = a_1 \alpha / 2$ and $\beta_2 = 2 a_2 \alpha$.
\end{proof}

The edge connectivity of Theorem~\ref{thm:random-connectivity} and the degree bounds of Lemma~\ref{lem:degree-bounds} bound the values $c(u,v)$.

\begin{lemma}
\label{lem:connectivity-bounds}
Let $G(n, p)$ be an Erdős–Rényi random graph with $p = \alpha \log n / n$, where $\alpha > 1$, and let $\beta_1, \beta_2$ be the constants from Lemma~\ref{lem:degree-bounds}.
Then
\[
\lim_{n\to\infty} \mathbb{P}\Bigl( \forall u \neq v: \ \beta_1 \log n \leq c(u,v) \leq \beta_2 \log n \Bigr) = 1.
\]
\end{lemma}

\begin{proof}
By Theorem~\ref{thm:random-connectivity} and Lemma~\ref{lem:degree-bounds}, with high probability, the graph $G(n, p)$ is $k$-edge-connected for $k = \delta(G)$ and all its degrees lie between $\beta_1 \log n$ and $\beta_2 \log n$.
Suppose that both events hold.
Since $G$ is $\delta(G)$-edge-connected, every cut separating two distinct vertices has at least $\delta(G)$ edges.
Hence all pairs of distinct vertices $u, v$ satisfy
\[
c(u, v) \geq \min_{x \neq y} c(x, y) \geq \min_{z \in V(G)} \deg(z) \geq \beta_1 \log n.
\]

On the other hand, the edges incident to $u$ form a cut that separates $u$ from $v$.
A minimum $u$--$v$ cut has at most as many edges as this cut, so
\[
c(u,v)\leq \deg(u) \leq \beta_2 \log n.
\]

Hence
\[
\beta_1 \log n \leq c(u,v) \leq \beta_2 \log n. \qedhere
\]
\end{proof}

By Lemma~\ref{lem:basic-bounds} and its proof, the union of a minimum $u$--$v$ cut and a shortest $u$--$v$ path has at most $c(u,v) + d(u,v) - 1$ edges, and every cut-path has at least $c(u,v)$ edges.
The diameter bound of Theorem~\ref{thm:random-diameter} and the connectivity bound of Lemma~\ref{lem:connectivity-bounds} keep the ratio of these two numbers below $1+\epsilon$ on almost all inputs.

\begin{theorem}
\label{thm:approximation-scheme}
Let $G(n, p)$ be an Erdős–Rényi random graph with $p \geq \alpha \log n / n$, where $\alpha > 1$.
Then, for every $\epsilon > 0$, the algorithm that returns the union of a minimum $u$--$v$ cut and a shortest $u$--$v$ path is an average $(1+\epsilon)$-approximation scheme for the \MinCutPath{} problem.
\end{theorem}

\begin{proof}
The algorithm runs in polynomial time, since both a minimum cut and a shortest path can be computed in polynomial time~\cite{Cormen2022}.
By the proof of Lemma~\ref{lem:basic-bounds}, its output $S$ is a cut-path with $|S| \leq c(u,v) + d(u,v) - 1$.

Adding edges to a graph cannot increase the distances and cannot decrease the values $c(u,v)$.
Hence, by the standard coupling of the random graphs $G(n,p)$ for different edge probabilities (see, e.g.,~\cite{Frieze2016}), the bounds of Theorem~\ref{thm:random-diameter} and Lemma~\ref{lem:connectivity-bounds}, which are stated for $p = \alpha \log n / n$, remain valid for every $p \geq \alpha \log n / n$.
Thus, with high probability, all pairs of distinct vertices $u,v$ satisfy
\[
d(u,v) \leq \diam(G) \leq \frac{\log n}{\log \log n}
\qquad\text{and}\qquad
c(u,v) \geq \beta_1 \log n.
\]
By Lemma~\ref{lem:basic-bounds}, the minimum cut is no larger than the minimum cut-path:
\[
c(u,v) \leq \cp(u,v) = \OPT(x).
\]

Hence
\begin{align*}
\frac{|S|}{\OPT(x)} &\leq \frac{c(u,v)+d(u,v)-1}{c(u,v)} \leq \frac{c(u,v)+d(u,v)}{c(u,v)} = 1 + \frac{d(u,v)}{c(u,v)}\\
&\leq 1+\frac{\log n / \log \log n}{\beta_1 \log n} = 1 + \frac{1}{\beta_1 \log \log n}.
\end{align*}
The right-hand side is smaller than $1+\epsilon$ for all sufficiently large $n$.
Thus, the algorithm achieves the approximation ratio $1+\epsilon$ on a set of inputs whose probability tends to one, i.e., it is an average $(1+\epsilon)$-approximation scheme.
\end{proof}

\section{Conclusion}
\label{sec:conclusion}

We introduced the \MinCutPath{} problem and proved that its decision version is NP-complete.
In graphs of diameter two and in graphs with minimum cut-value at most two, $\cp(u, v) = c(u, v) + d(u, v) - 1$, and the union of a minimum cut and a shortest path is a minimum cut-path.
Dense Erdős–Rényi random graphs have diameter two with high probability, and in sparser ones with $p \geq \alpha \log n / n$, $\alpha > 1$, the same union is within a factor of $1+\epsilon$ of the optimum on almost all inputs.

Several questions remain open.
First, in which further graph classes can the problem be solved in polynomial time?
A natural candidate is the class of graphs in which every pair of vertices is at distance at most two or has a minimum cut of size at most two; it contains both classes treated in this paper.
Is \MinCutPath{} NP-hard on graphs of diameter three, the smallest diameter for which finding the most vital edges for shortest paths and finding a matching cut are NP-hard~\cite{Bazgan2019MostVital,LeLe2019MatchingCut}?

Second, does duality help to solve \MinCutPath{} on \emph{planar graphs}?
In a planar graph, every minimal edge cut corresponds to a cycle in the dual graph and vice versa, and several problems on cuts and flows can be solved faster in planar graphs than in general ones~\cite{itai1979maximum,henzinger1997faster}.

Third, how do algorithms for \MinCutPath{} behave on average?
This work is purely theoretical.
Experiments on real-world networks and on random instances would answer this question and could guide the design of heuristics.

Finally, what is the structure of the instances $(G, u, v)$ in which $\cp(u,v)$ equals a fixed constant?
A closer study of them may clarify the boundary between polynomially solvable and NP-hard instances.


\section*{Acknowledgments}

The author thanks Martin Loebl, who proposed the \MinCutPath{} problem, for his guidance and for valuable discussions.

